\documentclass{article}
\usepackage{bbold}
\usepackage{amsfonts}
\usepackage{amsmath}
\usepackage{amssymb}
\usepackage{color}
\setlength{\columnseprule}{1pt}
\usepackage{cmap}
\usepackage[utf8]{inputenc}
\usepackage[T2A]{fontenc}
\usepackage[english, russian]{babel}
\usepackage{graphicx}
\usepackage{hyperref}
\usepackage{mathdots}
\usepackage{xfrac}


\def\columnseprulecolor{\color{black}}

\graphicspath{ {./resources/} }


\usepackage{listings}
\usepackage{xcolor}
\definecolor{codegreen}{rgb}{0,0.6,0}
\definecolor{codegray}{rgb}{0.5,0.5,0.5}
\definecolor{codepurple}{rgb}{0.58,0,0.82}
\definecolor{backcolour}{rgb}{0.95,0.95,0.92}
\lstdefinestyle{mystyle}{
    backgroundcolor=\color{backcolour},   
    commentstyle=\color{codegreen},
    keywordstyle=\color{magenta},
    numberstyle=\tiny\color{codegray},
    stringstyle=\color{codepurple},
    basicstyle=\ttfamily\footnotesize,
    breakatwhitespace=false,         
    breaklines=true,                 
    captionpos=b,                    
    keepspaces=true,                 
    numbers=left,                    
    numbersep=5pt,                  
    showspaces=false,                
    showstringspaces=false,
    showtabs=false,                  
    tabsize=2
}

\lstset{extendedchars=\true}
\lstset{style=mystyle}

\newcommand\0{\mathbb{0}}
\newcommand{\eps}{\varepsilon}
\newcommand\overdot{\overset{\bullet}}
\DeclareMathOperator{\sign}{sign}
\DeclareMathOperator{\re}{Re}
\DeclareMathOperator{\im}{Im}
\DeclareMathOperator{\Arg}{Arg}
\DeclareMathOperator{\const}{const}
\DeclareMathOperator{\rg}{rg}
\DeclareMathOperator{\Span}{span}
\DeclareMathOperator{\alt}{alt}
\DeclareMathOperator{\Sim}{sim}
\DeclareMathOperator{\inv}{inv}
\DeclareMathOperator{\dist}{dist}
\DeclareMathOperator{\supp}{supp}
\DeclareMathOperator{\Var}{Var}
\DeclareMathOperator{\Cov}{Cov}
\DeclareMathOperator{\argmax}{argmax}
\DeclareMathOperator{\argmin}{argmin}
\DeclareMathOperator{\med}{med}
\newcommand\1{\mathbb{1}}
\newcommand\ul{\underline}
\newcommand{\ppart}[2]{\frac{\partial #1}{\partial #2}}
\renewcommand{\bf}{\textbf}
\renewcommand{\it}{\textit}
\newcommand\vect{\overrightarrow}
\newcommand{\nm}{\operatorname}
\DeclareMathOperator{\df}{d}
\DeclareMathOperator{\tr}{tr}
\newcommand{\bb}{\mathbb}
\newcommand{\lan}{\langle}
\newcommand{\ran}{\rangle}
\newcommand{\an}[2]{\lan #1, #2 \ran}
\newcommand{\fall}{\forall\,}
\newcommand{\ex}{\exists\,}
\newcommand{\lto}{\leftarrow}
\newcommand{\xlto}{\xleftarrow}
\newcommand{\rto}{\rightarrow}
\newcommand{\xrto}{\xrightarrow}
\newcommand{\uto}{\uparrow}
\newcommand{\dto}{\downarrow}
\newcommand{\lrto}{\leftrightarrow}
\newcommand{\llto}{\leftleftarrows}
\newcommand{\rrto}{\rightrightarrows}
\newcommand{\Lto}{\Leftarrow}
\newcommand{\Rto}{\Rightarrow}
\newcommand{\Uto}{\Uparrow}
\newcommand{\Dto}{\Downarrow}
\newcommand{\LRto}{\Leftrightarrow}
\newcommand{\Rset}{\bb{R}}
\newcommand{\Rex}{\overline{\bb{R}}}
\newcommand{\Cset}{\bb{C}}
\newcommand{\Nset}{\bb{N}}
\newcommand{\Qset}{\bb{Q}}
\newcommand{\Zset}{\bb{Z}}
\newcommand{\Bset}{\bb{B}}
\renewcommand{\ker}{\nm{Ker}}
\renewcommand{\span}{\nm{span}}
\newcommand{\Def}{\nm{def}}
\newcommand{\mc}{\mathcal}
\newcommand{\mcA}{\mc{A}}
\newcommand{\mcB}{\mc{B}}
\newcommand{\mcC}{\mc{C}}
\newcommand{\mcD}{\mc{D}}
\newcommand{\mcJ}{\mc{J}}
\newcommand{\mcT}{\mc{T}}
\newcommand{\us}{\underset}
\newcommand{\os}{\overset}
\newcommand{\ol}{\overline}
\newcommand{\ot}{\widetilde}
\newcommand{\vl}{\Biggr|}
\newcommand{\ub}[2]{\underbrace{#2}_{#1}}
\newcommand{\ob}[2]{\overbrace{#2}^{#1}}
\newcommand{\pat}{\partial}

\def\letus{%
    \mathord{\setbox0=\hbox{$\exists$}%
             \hbox{\kern 0.125\wd0%
                   \vbox to \ht0{%
                      \hrule width 0.75\wd0%
                      \vfill%
                      \hrule width 0.75\wd0}%
                   \vrule height \ht0%
                   \kern 0.125\wd0}%
           }%
}
\DeclareMathOperator*\dlim{\underline{lim}}
\DeclareMathOperator*\ulim{\overline{lim}}

\everymath{\displaystyle}

% Grath
\usepackage{tikz}
\usetikzlibrary{positioning}
\usetikzlibrary{decorations.pathmorphing}
\tikzset{snake/.style={decorate, decoration=snake}}
\tikzset{node/.style={circle, draw=black!60, fill=white!5, very thick, minimum size=7mm}}
\newcommand{\norm}[1]{\left\lVert#1\right\rVert}
\title{\hugeТеория кодирования}
\author{Матвеев Сергей M3338}
\date{7 семестр}
\begin{document}

\maketitle

\section{Структура систем передачи информации. Простейшие методы модуляции. Кри-
терии идеального наблюдателя и максимума правдоподобия. Вероятность ошиб-
ки сигналов 2-АМ в случае канала с АБГШ. Отношение сигнал/шум на бит и
на символ.}


\section{Блоковые коды и их параметры. Критерии декодирования и метрики. Границы
Хемминга и Варшамова-Гилберта.}

\section{Линейные коды. Границы Синглтона, Варшамова-Гилберта и Грайсмера. Веро-
ятность ошибки декодирования и необнаружения ошибки.}

\section{Декодирование линейных кодов. Таблица стандартной расстановки. Код Хем-
минга. Стирания.}

\section{Декодирование по информационным совокупностям.}

\section{Дуальные коды. Весовой спектр кода. Тождество Мак-Вильямс и его доказательство.}

\section{Критерии мягкого декодирования. Метод порядковых статистик.}

\section{Минимальная решетка линейного блокового кода и способы ее построения. Алгоритм Витерби.}

\section{Декодирование с мягким выходом. Алгоритм БКЕР.}

\section{Сверточные коды. Способы представления. Катастрофические порождающие матрицы.}

\section{Производящая функция сверточных кодов. Вероятность ошибки декодирования сверточных кодов с помощью алгоритма Витерби.}

\section{Конструкция Плоткина. Коды Рида-Маллера. Прямое произведение кодов. Обобщенные каскадные коды.}

\section{Укорочение, выкалывание и расширение кодов. Турбо-коды.}

\section{Параметр Бхаттачарьи. Поляризация канала. Полярные коды. Сложность кодирования.}

\section{Алгоритм последовательного исключения и декодер Тала-Варди.}

\section{Построение полярных кодов. Полярные коды с CRC, полярные подкоды.}

\section{Циклические коды. Порождающий и проверочный многочлены. Кодирование.}

\section{Идеалы, факторкольца и поля.}

\section{Конечные поля. Характеристика поля. Число элементов в поле Галуа и его примитивный элемент. Основные свойства.}

\section{Минимальные многочлены.}

\section{Проверочная матрица циклического кода над расширенным полем. Коды БЧХ. Граница БЧХ. Свойства.}

\section{Коды Рида-Соломона и обобщенные коды Рида-Соломона.}

\section{Алгоритм Питерсона-Горенстайна-Цирлера декодирования кодов БЧХ.}

\section{Ключевое уравнение декодирования кодов БЧХ. Алгоритм Форни. Декодирование с помощью алгоритма Евклида.}

\section{Альтернантные коды. Коды Гоппы. Криптосистема Мак-Элиса.}

\section{Низкоплотностные коды. Основные характеристики. Конструкции низкоплотностных кодов.}

\section{Декодирование низкоплотностных кодов.}

\section{Эволюция плотностей и порог итеративного декодирования низкоплотностных кодов.}

\section{Кодирование в стирающих каналах. Цифровой фонтан. LT-коды и хищные коды.}

\section{Битоперемеженная кодовая модуляция. Многоуровневые коды.}


\end{document}
